\centerline{\bf \Huge Инверсия. Начало}

\begin{flushright}
        \scriptsize{Эпизод 12. Ценность чуда}
\end{flushright}

\problem{1}
Дана окружность $\omega$ с центром в точке $O$ и точка $P$ на этой окружности. Докажите, что после инверсии с центром в точке $P$ окружность $\omega$ перейдет в некоторую прямую $\ell$, а точка $O$ перейдет в точку, симметричную $P$ относительно $\ell$.

\problem{2}
(a) Окружности $\omega_1, \omega_2, \omega_3$ и $\omega_4$ касаются друг друга по циклу. Докажите, что четыре точки касания лежат на одной окружности.

\noindent
(b) Окружности $\omega_1, \omega_2, \omega_3$ и $\omega_4$ пересекают друг друга по циклу в двух точках: красной и синей. Оказалось, что четыре красные точки лежат на одной окружности. Докажите, что четыре синие точки тогда также лежат на одной окружности.

\problem{3}
(a) Дан треугольник $A B C$. Обозначим за $\Gamma$ окружность, проходящую через вершины $B$ и $C$, а также центр вписанной окружности этого треугольника. Куда переходит ( $A B C$ ) при инверсии относительно $\Gamma$ ?

\noindent
(b) На биссектрисе угла $A$ треугольника $A B C$ выбираются изогонально сопряженные точки $P$ и $Q$. Окружность $\omega$ проходит через точки $P, Q$ и касается отрезка $B C$. Докажите, что $\omega$ и $(A B C)$ касаются.

\problem{4}
\textbf{Неравенство Птолемея.}
Для четырех точек $A, B, C$ и $D$ на плоскости докажите неравенство $A B \cdot C D+B C \cdot A D \geqslant A C \cdot B D$. Когда достигается равенство?

\problem{5}
Дан полукруг $\Gamma$ с диаметром $A B$. Точка $C$ выбирается на отрезке $A B$, после чего на отрезках $A C$ и $B C$ строятся два полукруга $\gamma_a$ и $\gamma_b$ внутрь Г. Перпендикуляр, восставленный к прямой $A B$ в точке $C$, пересекает отрезок $\Gamma$ в точке $D$. Рассматриваются криволинейные треугольники: первый образован дугами $\Gamma, \gamma_a$ и отрезком $C D$, а второй - дугами $\Gamma, \gamma_b$ и отрезком $C D$. Докажите, что вписанные окружности этих треугольников равны.

\problem{6}
Через точку $A$ вне окружности с центром в $O$ проведены касательные $A X$ и $A Y$ к этой окружности, а также секущая, пересекающая окружность в точках $Z$ и $T$. Докажите, что точки $Z$, $T, O$ и середина $X Y$ лежат на одной окружности.

\problem{7}
На плоскости провели несколько окружностей и отметили все точки их пересечения или касания. Может ли оказаться, что на каждой окружности лежат ровно пять отмеченных точек, а через каждую отмеченную точку проходят ровно пять окружностей?

\newpage

\hspace{3mm}

\problem{8}
В треугольнике $A B C$ провели окружности $\omega_a, \omega_b$ и $\omega_c$, вписанные в углы $A, B$ и $C$ и проходящие через инцентр $I$ этого треугольника. Обозначим за $A_1, B_1$ и $C_1$ вторые пересечения этих трех окружностей друг с другом. Докажите, что центры окружностей $A A_1 I, B B_1 I$ и $C C_1 I$ лежат на одной прямой.